\documentclass[12pt,letterpaper]{article}
\usepackage[T1]{fontenc}
\usepackage{lmodern,microtype,amsmath,amssymb,amsthm,xcolor}
\usepackage[margin=1.04in]{geometry}
\usepackage[colorlinks=true,linkcolor=black,citecolor=black,urlcolor=blue!45!black]{hyperref}
\hypersetup{pdftitle={Almost-everywhere circuit lower bounds from sparse 3-SAT}}
\newtheorem{theorem}{Theorem}
\newtheorem{lemma}[theorem]{Lemma}
\newcommand{\ENP}{\mathrm{E}^{\mathrm{NP}}}
\newcommand{\size}{C_{B_2}}
\newcommand{\poly}{\operatorname{poly}}
\newcommand{\bits}{\{0,1\}}
\title{Almost-everywhere circuit lower bounds\\from sparse 3-SAT}
\author{}
\date{23 November 2026}
\begin{document}
\maketitle
\begin{abstract}
Suppose 3-SAT on \(V\) variables and at most \(5V\) clauses has a deterministic
algorithm running in \(V^{O(1)}2^{O(V/\sqrt{\log V})}\) time. Then one language
in \(\ENP\) requires \(\Omega(n\sqrt{\log n})\) unrestricted fan-in-two gates
at every sufficiently large input length \(n\). We adapt the bounded-witness
refuter of Chen, Lyu and Williams to select a hard accepting computation at
each address length. The verifier reads the occupied portion of the initial
tape directly, so the number of variables in its SAT query grows linearly
with the gate budget and the address length. The same construction transfers
general nondecreasing savings rates without having to compute the rate.
\end{abstract}

\section{The transfer}
Write \(W_n\) for the characteristic function of a language on \(\bits^n\).
The size \(\size(f)\) counts noninput gates in a single-output circuit over
all binary Boolean functions, with arbitrary depth and free fan-out.
All logarithms have base two. We use the multitape bit model, and \(\ENP\)
allows a deterministic \(2^{O(n)}\)-time machine to write queries of that length
to an NP oracle.

\begin{theorem}\label{thm:transfer}
Let \(a(V)\ge1\) be nondecreasing and tend to infinity. Suppose a uniform
deterministic algorithm \(D\) decides satisfiability for CNFs of width at
most three on \(V\) declared variables with at most \(5V\) clauses in time
\(O(V^b 2^{A V/a(V)})\), for fixed \(A,b>0\). The declaration may include
unused variables. There exist \(W\in\ENP\) and \(c,\rho>0\) such that for
every sufficiently large \(n\),
\[
       \size(W_n)>c\min\{n\,a(n),\,2^{\rho n}\}.
\]
The language can be chosen as a function of \(D\) alone.
\end{theorem}

In particular, if \(\log a(n)=o(n)\), the conclusion is
\(\size(W_n)>c\,n a(n)\) eventually. At the fixed-density rate reported in
October 2026 this gives \(c\,n\sqrt{\log n}\). OpenAI's Corollary 1.4
obtains this size at infinitely many lengths by applying the nondeterministic
time hierarchy to a unary language \cite{oai}. Here the hard function is a
tableau chosen separately at each length.

Fortnow and Santhanam proved an almost-everywhere hierarchy for machines with
bounded nondeterminism \cite{fs}. Chen, Lyu and Williams gave a refuter that
finds a disagreement at a prescribed length \cite[Secs.~4--5]{clw}. Their
Theorem 1.1 allows a fixed power of the target size in the hypothesis on
Gap-UNSAT. At the present gate scale that power matters. We use their
root-and-successor computation and account directly for the size of its local
test.

\section{One tape and a small test for a guessed tableau}
Fix a binary-branching, one-tape interpreter \(U\) with a left end. Its input
\(x(e,n,s,w)\) encodes a deterministic multitape verifier \(e\), unary \(n\),
a positive integer gate budget \(s\) in binary, and a raw bit string \(w\).
Give \(e\) and the binary expansion of \(s\) a unary length prefix followed
by a zero. A zero terminates unary \(n\), and the first blank terminates \(w\).
Write \(g=|w|\). The two inputs to a simulated verifier are a word and a
separate witness.

First \(U\) runs \(e\) on the \emph{root} \(x(e,n,s,0^g)\), supplying \(w\)
as its witness. If that branch accepts, \(U\) rejects. If it rejects and
\(w=1^g\), \(U\) accepts. Otherwise, after root rejection, \(U\) runs \(e\)
on \(x(e,n,s,w+1)\) with a freshly guessed \(g\)-bit witness and adopts its
answer. Here \(w+1\) is the fixed-width binary successor. Acceptance returns
the head to the origin, writes a distinguished symbol and enters an absorbing
state. Rejection is absorbing as well.

For each fixed \(e\), if both simulated branches take at most \(t\) steps
and \(|x|=m\), the whole computation takes at most \(C_e(t+m+1)^3\) steps.
Indeed the interpreter stores the tapes as marked strings. The total used
length is \(O_e(t+m+1)\). Sweeps collect the scanned symbols, find the
instruction in the stored program, and change symbols and head marks. Even a
quadratic allowance in the used length per simulated step suffices. Copying
the root, forming the successor, preparing either witness and returning to
the origin fit the same bound. Only the multiplicative constant depends on
the simulated program.

Put \(r=\lfloor n/5\rfloor\), \(R=2^r\). A table on \(\bits^n\) describes
\(R\) rows of \(2R\) cells. Use a tag, \(r\) row bits, \(r+1\) column bits
and a fixed number of component-selector bits for a cell's symbol and
head state. Set the remaining address bits to zero. A second tag with the
column and selector zero specifies one choice bit for each transition row.
The used fields fit at all large \(n\). Other table bits are unrestricted.

The first row contains \(x\), followed by blanks, with one head at the origin.
The last row must carry the distinguished accepting head and symbol there.
For each transition, the old three-cell neighborhood and the choice bit
determine the new center cell. Use outside sentinels at the two ends, reject
outward moves, and test legal cell encodings. The address for the choice
depends only on the old row. Starting with the unique initial head, induction
therefore gives precisely the genuine computation for the row choices.
Conversely, if \(m<R\), any accepting branch of fewer than \(R-1\) steps
fits in the rectangle and extends to its final row by absorption.

\begin{lemma}\label{lem:checker}
For \(m<R\) and address fields that fit, given \(x\) and an \(n\)-input circuit
of at most \(s\) gates, validity of its table can be decided by polynomial-time
preprocessing and at most one call to \(D\) with \(V=K(n+s)\) variables.
Here \(K\) is a fixed integer depending only on \(U\).
\end{lemma}
\begin{proof}
For each \(j<m\), evaluate the supplied circuit at all component addresses
of cell \(j\) in row zero. Compare the resulting cell to the symbol \(x_j\),
including the head and the left-end convention. There are only
\(O_U(m)\) evaluations, each taking \(\poly(n+s)\) bit operations by
scanning the gate list.

For the remaining checks build a circuit whose inputs select a test, row
and column, and whose output indicates a violation. The blank portion of
row zero uses the guard \(j\ge m\). Transition and final-state tests use
only fixed predicates. A constant number of copies of the guessed circuit
read the cell components and the common choice bit. Comparing indices,
testing endpoints and shifting a row or column by one take \(O(n)\) gates.
This gives \(G=O_U(s+n)\) gates and \(O(n)\) inputs, even when \(m\) is
large. Primary inputs and gate outputs receive distinct variables.
For each binary gate, four clauses forbid its inconsistent truth-table
rows. Pin the violation output to one with a unit clause. Each primary
assignment has a unique consistent extension to gate outputs. The circuit
has a violation exactly when this formula is satisfiable.
Declare unused variables up to \(V=K(n+s)\), choosing \(K\ge1\) large enough
to cover all variables and \(4G+1\le 5V\). Accept exactly when the occupied
cells passed and the formula is unsatisfiable.
\end{proof}

\section{Choosing a gate budget under one clock}
Let \(M\) be the verifier in Lemma~\ref{lem:checker}, using the actual input
word \(x(e,n,s,w)\) as the initial tape to check. Its separate witness \(u\)
describes an \(s\)-gate circuit. With \(\ell=\lceil\log(n+s+2)\rceil\), use
\[
          g=g(n,s)=s(4+2\ell)+\ell
\]
bits for the gate truth tables, two predecessor indices per gate, and the
output index. Reject malformed words, wrong lengths, forward references
and out-of-range indices. Unused dummy gates allow every circuit of size
at most \(s\) to be encoded. Fix \(e=\operatorname{code}(M)\) now. The
alphabet and transition rule of \(U\), hence \(K\), were fixed before
specifying \(M\). Constructing the formula uses the transition rule as a
finite table.

For chain inputs \(|w|=|u|=g\) the length \(m\) depends only on \(e,n,s\)
and is \(O_e(n+(s+1)\log(n+s+2))\). Whenever the fields fit and \(m<R\),
the time on \emph{every} pair \(w,u\) has logarithm at most
\begin{equation}
     A\,\frac{K(n+s)}{a(K(n+s))}+B\log(n+s+2)+O_e(1) . \label{eq:cost}
\end{equation}
Here \(B\) absorbs the solver's polynomial factor, the scans over the input
and circuit, and all formula construction. The verifier need not know \(a\).

Set \(t=2^{\lfloor n/40\rfloor}\). Call a budget \emph{admissible} at \(n\)
if the address fields fit, \(m\le t<R\), and \(M\) halts in at most \(t\)
steps on all \(g\)-bit pairs \(w,u\) at that budget. A deterministic
\(\ENP\) machine can find the largest admissible budget among
\(1,2,4,\ldots,2^n\), using zero output if there is none. First discard
budgets failing the size checks. For each remaining budget ask if there
\emph{exists} a pair \(w,u\) on which \(M\) is still running after \(t\)
steps. With \(t\) written in unary, an NP machine guesses the pair and
simulates the deterministic verifier. A negative answer certifies the
budget. There are only \(n+1\) budgets, and each gets its own test.

For comparison, choose small constants \(0<\eta,\rho<1\). Take \(s\) to
be the largest power of two at most \(\min\{\eta n a(n),2^{\rho n}\}\).
Since \(K(n+s)\ge n\), monotonicity bounds the first term of
\eqref{eq:cost} by \(AK\eta n+AKn/a(n)\). Also
\(\log(n+s+2)\le\rho n+o(n)\) and \(\log m\le\rho n+o(n)\).
Choose \(AK\eta+B\rho<1/80\) and \(\rho<1/80\).
Then \(m\le t\) and every branch halts by \(t\) for all large \(n\).
The largest admissible budget is at least this comparison budget, hence
at least \((\eta/2)\min\{n a(n),2^{\rho n}\}\).
Only the proof chooses \(\eta,\rho,s\) from the rate.

\paragraph{The square-root logarithm case.}
For this rate one may fix a small positive rational \(\eta\), set
\(s=\lfloor\eta n\lfloor\sqrt{\lfloor\log n\rfloor}
\rfloor\rfloor\) and omit the budget search. Estimate~\eqref{eq:cost} gives
log-time at most \((AK\eta+o(1))n\). A lookup circuit in the style of
Lupanov \cite{lupanov} can also check the whole initial row inside the SAT
query. At this budget \(m=O(s\log n+n)\).
Pad the \(m\)-bit word to a power of two and use blocks of power-of-two length
\(b\in[\tfrac14\log m,\tfrac12\log m]\). Prefix trees produce the block
indicators and the within-block indicators in \(O(m/b+b)\) gates. Form all
\(2^b\) block words as disjunctions of the latter indicators, then use each
block indicator to select its prescribed word. The lookup uses
\(O(m/b+b2^b)=O(m/\log m)\) gates and is constructible in polynomial time.
Since \(m\ge n\), this adds \(O(s+n)\) gates with a constant independent of
\(\eta\le1\). A comparator supplies blanks beyond the word.

\section{Finding a hard table at the chosen length}
Fix a sufficiently large \(n\) and its selected admissible \(s\).
The interpreter uses at most \(C_e(2t+1)^3<R-1\) steps on each branch
of every chain input. List the \(g\)-bit strings as
\(w_0,\ldots,w_{H-1}\) in binary order, with \(H=2^g\). Write
\(x_i=x(e,n,s,w_i)\). Consider three bits for each \(i\).
\begin{align*}
 P_i & = [M(x_i,u)\text{ accepts for some }u\in\bits^g],\\
 d_i & = [M(x_0,w_i)\text{ accepts}],\\
 Q_i & = [U(x_i)\text{ has an accepting table}].
\end{align*}
Set \(P_H=1\). The root branches enumerate all witnesses, the two simulations
fit in the rectangle, and every accepted circuit passes the exact table test.
Consequently,
\[
       P_0=\bigvee_{i<H}d_i,\qquad
       Q_i=(1-d_i)P_{i+1},\qquad P_i\le Q_i \quad(0\le i<H).
\]
If \(P_0=1\), induction gives \(P_i=1\) and \(d_i=0\) at each actual
index, contradicting the disjunction. Thus \(P_0=0\) and every \(d_i=0\).
Between the zero at \(0\) and the sentinel one at \(H\), retain a zero
left endpoint and a one right endpoint by querying their midpoint. After
at most \(g\) queries this gives consecutive endpoints \(j,j+1\).
Now \(P_j=0\) and \(Q_j=P_{j+1}=1\). There is an accepting table for
\(x_j\), and every circuit of at most \(s\) gates fails to encode one.
Indeed, a circuit for any such table would supply an accepting description
to \(M(x_j,\cdot)\).

On input \(y\in\bits^n\), compute that same \(j\), using an NP query with
unary clock \(t\) for each \(P_i\). The sentinel is supplied directly.
The endpoints need \(g+1\) bits, and \(g\le m\le t\). Write out the
ordinary initial, transition and final constraints on \(2^n\) table
bits for \(x_j\). Their total encoding length is \(2^{O(n)}\).
Select the lexicographically first satisfying table, in address order,
by testing whether each successive zero bit has a satisfying extension
of the chosen prefix. Return its bit at \(y\). Use zero if the formula
is unsatisfiable or no budget was selected. At most \(2^n+1\) SAT queries
of length \(2^{O(n)}\) suffice for the table selection.
The algorithm therefore defines one \(W\in\ENP\). At each sufficiently
large \(n\) its slice encodes an accepting table on \(x_j\), so
\(\size(W_n)>s\). This proves Theorem~\ref{thm:transfer}.

\section{Instantiating the decision bound}
OpenAI's proof of Corollary 1.4 explicitly uses a decision algorithm for
\(V\) variables and at most \(5V\) clauses in time
\(V^{O(1)}2^{O(V/\sqrt{\log V})}\). The polynomial-space implementation
allows a multitape simulator to scan the workspace on each bit access
with polynomial overhead per access \cite{oai}. OpenAI's separate Lean
formalization records the exponent \(O(V(\log\log V)^2/\sqrt{\log V})\).
Using that bound alone gives \(\Omega(n\sqrt{\log n}/(\log\log n)^2)\)
at every sufficiently large length by the same theorem.

Anthropic's \texttt{MainTheorem} gives a second route \cite{anthropic}.
For width three, the clause-count term has exponent \(O(q/\sqrt{\log q})\)
for \(q\ge2\) distinct clauses and at least four occurring variables.
Delete tautologies and duplicate clauses, simplify repeated literals,
drop unused variables and renumber the rest. An empty clause decides
unsatisfiability immediately. If fewer than four variables or fewer than
two clauses remain, a constant-size test suffices. Otherwise encode the
normalized formula in polynomial time, run the finite-stack machine and
test whether the returned count is positive. Each stack of a fixed
finite-alphabet machine can be maintained on its own tape. Since
\(q\le5V\) and \(u/\sqrt{\log u}\) increases for real \(u\ge2\), the
exponent is \(O(V/\sqrt{\log V})\). Either stated sparse bound supplies
the same almost-everywhere corollary.

\begin{thebibliography}{9}\small
\bibitem{lupanov} O. B. Lupanov. On a method of circuit synthesis.
\emph{Izvestiya VUZ, Radiofizika} 1(1), 120--140, 1958.
\href{https://radiophysics.unn.ru/issues/1958/1/120}{Original article}.
\bibitem{fs} L. Fortnow and R. Santhanam.
New non-uniform lower bounds for uniform classes. CCC, 2016.
\href{https://doi.org/10.4230/LIPIcs.CCC.2016.19}{Proceedings}.
\bibitem{clw} L. Chen, X. Lyu and R. Williams.
Almost-everywhere circuit lower bounds from non-trivial derandomization.
FOCS, 2020. \href{https://eccc.weizmann.ac.il/report/2020/150/}{Full version, ECCC TR20-150}.
\bibitem{oai} OpenAI. The Exponential Time Hypothesis is False.
Revised 25 October 2026. Theorem 1.1 and Corollary 1.4.
\href{https://github.com/openai/math/tree/main/preprints/The-Exponential-Time-Hypothesis-is-False-October-25-2026}{Manuscript},
\href{https://raw.githubusercontent.com/openai/math/7ef5c0cc5ca4b0e072b7a050fca09ab24bc6fa3a/preprints/The-Exponential-Time-Hypothesis-is-False-October-25-2026/build/quantitative-circuits.tex}{quantitative circuit proof}
and \href{https://github.com/openai/math/blob/main/lean/docs/378.md}{scope of the Lean formalization}.
\bibitem{anthropic} Claude (Anthropic). A \(2^{o(n)}\) Time Algorithm for 3-SAT.
October 2026. \href{https://github.com/anthropics/formal-math/blob/16a5ce41304fcdaea4e0be7c132e1c361dc26da5/eth/MainResults.lean}{Formal statement}
and \href{https://github.com/anthropics/formal-math/blob/16a5ce41304fcdaea4e0be7c132e1c361dc26da5/eth/Solution.lean}{proof endpoint}.
\end{thebibliography}
\end{document}
